\subsubsection{毛管波制約の導出：分散関係から CFL 条件へ}
\label{app:capillary_cfl}
% ═══════════════════════════════════════════════════════════════════════════════

\textbf{Step 1：毛管波の分散関係}

気液界面上の微小振幅波（毛管波）の線形安定解析から，
角周波数 $\omega$ と波数 $k$ の関係（分散関係）は：
\begin{equation}
  \omega^2 = \frac{\sigma\,k^3}{\rho_l + \rho_g}
  \label{eq:capillary_dispersion_app}
\end{equation}
位相速度 $c_\sigma = \omega/k = \sqrt{\sigma k/(\rho_l+\rho_g)}$，
群速度 $c_g = d\omega/dk = (3/2)c_\sigma$．

\textbf{Step 2：格子上の最大分解可能波数}

格子間隔 $\Delta x$ の離散格子で分解できる最高波数はナイキスト波数
$k_{\max} = \pi/\Delta x$ である．

\textbf{Step 3：CFL 条件の適用}

$k_{\max}$ に対応する群速度で CFL 条件 $c_g(k_{\max})\,\Delta t \le \Delta x$ を適用する：
\begin{align*}
  c_g(k_{\max}) &= \frac{3}{2}\sqrt{\frac{\sigma\pi/\Delta x}{\rho_l+\rho_g}} \\
  \Delta t &\le \frac{\Delta x}{c_g(k_{\max})}
  = \frac{2}{3}\sqrt{\frac{(\rho_l+\rho_g)\,\Delta x^3}{\sigma\pi}}
\end{align*}
厳密な係数は $2/(3\sqrt{\pi}) \approx 0.376$ であるが，
文献で広く用いられる簡略形 $1/\sqrt{2\pi} \approx 0.399$ を採用すると
式 \eqref{eq:dt_sigma} を得る（$6\%$ 大きい $\Delta t$ を許すが，
毛管波 CFL の安全係数 $C_\sigma \in (0,1]$ を
$\Delta t \le C_\sigma\Delta t_\sigma$ として掛けるため実用上問題ない）：
\[
  \Delta t_\sigma = \sqrt{\frac{(\rho_l+\rho_g)\,\Delta x^3}{2\pi\sigma}}
\]

\textbf{物理的解釈：} $\Delta t_\sigma \propto \Delta x^{3/2}$ は分散関係 $\omega \propto k^{3/2}$ に由来する．
格子を半分にすると $\Delta t_\sigma$ は $1/2^{3/2} \approx 0.354$ 倍（$\approx 2.8$ 倍制約強化）．

\medskip
\textbf{大密度比（$\rho_l \gg \rho_g$）での解釈：}
$\rho_l/\rho_g = 1000$（水/空気）では $\rho_l + \rho_g \approx \rho_l$（気相寄与は $0.1\%$ 以下）．
表面張力が\textbf{復元力}（ばね），液相 $\rho_l$ が\textbf{慣性}（質量）を担い，
気相の慣性は無視できる．密度比 1000 の系では密度比 1 の系より
$\sqrt{1000} \approx 32$ 倍大きな $\Delta t_\sigma$ が許される（制約緩和）．\qed

% ═══════════════════════════════════════════════════════════════════════════════
%% 旧 appendix_numerics_schemes_s4.tex（ALE 効果）を統合
% ═══════════════════════════════════════════════════════════════════════════════

\subsubsection{時間変化格子と ALE 効果の補足}
\label{app:grid_ale}

本稿の界面追随格子の標準定義は §\ref{sec:grid_ale} に置く．
本付録では，式~\eqref{eq:ale_omission_error} のスケール見積もりだけを補足する．

ALE 定式化では格子速度 $\bm{v}_\text{mesh}$ を輸送方程式に追加する必要がある：
$\partial\psi/\partial t + \nabla\cdot((\bm{u} - \bm{v}_\text{mesh})\psi) = 0$．
同じ相対速度 $\bm{u}-\bm{v}_\text{mesh}$ は運動量・密度輸送にも現れるため，
$\psi$ だけを補間する格子追従は閉じた保存系ではない．

\textbf{誤差の発生機構：}
1ステップあたりの界面移動量を $\delta = |\bm{u}|\Delta t$ とすると，
ALE 補正を省略することで $\bm{v}_\text{mesh}$ による追加対流量が無視される：
\[
  \varepsilon_{\text{ALE}} \sim |\bm{v}_\text{mesh}| \cdot |\bnabla\psi| \cdot \Delta t \sim U_\Gamma \cdot \frac{1}{\varepsilon} \cdot \Delta t
\]
ここで $|\bnabla\psi| \sim 1/(4\varepsilon)$（界面近傍プロファイル），$U_\Gamma$ は界面速度である．
Mode 1/Mode 2 の判定，および大密度比・高速界面で必須となる保存的再写像，
速度再投影，履歴リセット，面幾何再構築は §\ref{sec:grid_ale} の本編定義に従う．

% ═══════════════════════════════════════════════════════════════════════════════
